%% kapa: exam=1 round=34 year=2023 period=2 form=A subject=appraisal_law subject_ko=감정평가관계법규 seq=34 number=34
%% source: Q-Net 감정평가사 기출문제 (cst00309) artlSeq=5211789 file=raw/1008/5211789/2023년도 제34회 감정평가사 1차시험 2교시.pdf
%% conversion: kordoc 4.24.51 -> tools/gen1.py (auto); review=proofread-sonnet; flags=none
\begin{question}{34}
부동산등기법령상 이의에 관한 설명으로 옳은 것을 모두 고른 것은?
\begin{framedtext}
ㄱ. 새로운 사실이나 새로운 증거방법을 근거로 이의신청을 할 수 있다. \par
ㄴ. 등기관은 이의가 이유 없다고 인정하면 이의신청일부터 7일 이내에 의견을 붙여 이의신청서를 관할 지방법원에 보내야 한다. \par
ㄷ. 이의에는 집행정지의 효력이 없다. \par
ㄹ. 송달에 대하여는 「민사소송법」을 준용하고, 이의의 비용에 대하여는 「비송사건절차법」을 준용한다. \par
\end{framedtext}
\begin{choices}
  \item ㄱ, ㄴ
  \item ㄱ, ㄷ
  \item ㄱ, ㄹ
  \item ㄴ, ㄷ
  \item ㄷ, ㄹ
\end{choices}
\end{question}
\answer{5}
